% =====================================================================
% 数据说明与预处理（精简版）
% ---------------------------------------------------------------------
% 使用方式：替换主文档中的“数据说明与预处理”一节；依赖图片与主文档一致：
%   Q1_price_load_pv.pdf     典型日负载/光伏/电价
%   att2_energy_b.png        全年净负载持续曲线
%   att4_panel_b.png         全年电价日内分布（中位数+四分位带）
%   att2_seasonal_{winter,spring,summer,autumn}.png  四季日内曲线
% 精简说明：原 4 个独立 figure 合并为 1 个 2×2 组合图；四季四张小图用
%   tabular 内嵌排布，避免 subfigure 内 \\ 换行导致的文字重叠；原“全年
%   日电量”（att2_energy_a.png）与“全年电价日历热图”（att4_panel_a.png）
%   信息已折入正文，不再单独成图，以压缩 1–2 页。
% =====================================================================

\section{数据说明与预处理}

\subsection{数据来源与说明}

本题五份附件均为 Excel 格式：附件 1 是典型日分时电价、负载与光伏预测，支撑问题一的确定性调度；附件 2 是全年 365 天负载与光伏实测，用于问题二的滚动预测与回测；附件 3、4 分别是光伏日前预报与全年电价，用于问题三、四；附件 5 是储能参数与结果模板。字段口径见表~\ref{tab:fields}。

\begin{table}[htbp]
\centering
\small
\caption{数据字段说明}
\label{tab:fields}
\begin{tabular}{llll}
\toprule
字段名称 & 字段含义 & 数据类型 & 单位 \\
\midrule
电价 & 分时购电价格（附件1 单日、附件4 全年） & 数值型 & 元/kWh \\
小区负载 & 小区用电功率（附件1、附件2） & 数值型 & kW \\
光伏发电功率 & 光伏出力（附件1 预测、附件2 实测） & 数值型 & kW \\
光伏日前预报 & 按发布时刻给出的光伏预报（附件3） & 数值型 & kW \\
日期 / 时间 & 10 min 间隔的记录时刻 & 时间型 & — \\
\bottomrule
\end{tabular}
\end{table}

\subsection{数据预处理}

逐表核查：附件 1、2、4 的数值字段缺失值、负值、重复行均为 0，附件 2、4 日期连续 365 天；附件 3 数值字段无缺失、无负值，仅“日期”列因每天 4 次预报合并单元格，按前向填充补齐。按逐时段 IQR 检出的离群点（负载 1308、光伏 34、电价 89）逐一核对均落在冬季负载高峰、夏季强光、峰谷尖峰等真实波动上，故不做异常值剔除与标准化。真正需要处理的是时间口径、单位、夜间噪声与面向模型的特征构造：

\begin{enumerate}[label=\arabic*.)]
  \item \textbf{时间对齐}：附件时刻为时段起点（$0{:}00{+}1$ 表示次日 $00{:}00$--$00{:}10$），与结果模板表头一致；优化覆盖当天 $00{:}00$--$24{:}00$ 共 144 段，全年负载/光伏将前一日末段接到当天首段。
  \item \textbf{单位换算}：目标函数与储能递推均用电量（kWh），故按 $E=P\,\Delta t$（$\Delta t=1/6$\,h）把附件功率换算为电量。
  \item \textbf{光伏夜间置零}：全年恒为 0 的 43 个时段强制置 0；其余晨昏过渡段残值最大约 $10$\,kW，相对负载数千 kW 可忽略，保留不处理。
  \item \textbf{附件 3 尺度对齐}：预报按整点给出，与 10 min 实测尺度不一致，按被预报钟点提取整点值与实测对齐。
  \item \textbf{特征构造}：为 XGBoost 构造时刻编号、星期、月份、年内日序，以及前 $1/2/7$ 天同刻值、近 $3/7/30$ 天同刻均值、$7$ 天同刻标准差共 11 个特征，且只用预测日前数据，避免信息泄漏。
  \item \textbf{残差分位保守修正}：点预测偏差会转化为供电缺口，故问题二将净负载残差的 $\alpha=0.8$ 高分位叠加到点预测上（见问题二），使计划购电偏保守、降低紧急购电风险。
\end{enumerate}

\subsection{数据探索性分析}

核心指标描述性统计见表~\ref{tab:desc}。

\begin{table}[htbp]
\centering
\small
\caption{核心指标描述性统计}
\label{tab:desc}
\begin{tabular}{lccccc}
\toprule
指标 & 均值 & 标准差 & 最小值 & 最大值 & 中位数 \\
\midrule
电价（附件1，元/kWh） & 0.77 & 0.31 & 0.37 & 1.40 & 0.81 \\
小区负载（附件1，kW） & 4626.03 & 977.47 & 3309.39 & 5958.97 & 4377.57 \\
光伏预测（附件1，kW） & 2311.78 & 2854.40 & 0 & 7612.32 & 387.61 \\
小区负载（附件2 全年，kW） & 4626.03 & 1411.73 & 1995.72 & 7978.88 & 4270.68 \\
光伏实测（附件2 全年，kW） & 2311.79 & 2948.64 & 0 & 10216.20 & 18.86 \\
\bottomrule
\end{tabular}
\end{table}

由表~\ref{tab:desc} 可见：光伏全年中位数仅 $18.86$\,kW、均值却达 $2311.79$\,kW，二者相差悬殊，说明光伏存在大量夜间零出力时段、白天集中出力；负载全年极差 $1995.72\sim7978.88$\,kW、标准差 $1411.73$\,kW，波动明显大于典型日，须逐日滚动预测而非套用固定模板。图~\ref{fig:eda} 把典型日、全年净负载、四季曲线与全年电价整合为一张图。

\begin{figure}[htbp]
  \centering
  \begin{subfigure}[t]{0.48\textwidth}
    \centering
    \includegraphics[width=\linewidth]{Q1_price_load_pv.pdf}
    \caption{典型日负载、光伏与分时电价}
    \label{fig:eda_typical}
  \end{subfigure}\hfill
  \begin{subfigure}[t]{0.48\textwidth}
    \centering
    \includegraphics[width=\linewidth]{att2_energy_b.png}
    \caption{全年净负载持续曲线}
    \label{fig:eda_netload}
  \end{subfigure}

  \vspace{1.2em}

  \begin{subfigure}[t]{0.48\textwidth}
    \centering
    \includegraphics[width=\linewidth]{att4_panel_b.png}
    \caption{全年电价日内分布}
    \label{fig:eda_price}
  \end{subfigure}\hfill
  \begin{subfigure}[t]{0.48\textwidth}
    \centering
    \begin{tabular}{@{}c@{\hspace{0.6em}}c@{}}
      \includegraphics[width=0.22\textwidth]{att2_seasonal_winter.png} &
      \includegraphics[width=0.22\textwidth]{att2_seasonal_spring.png} \\[0.4em]
      \includegraphics[width=0.22\textwidth]{att2_seasonal_summer.png} &
      \includegraphics[width=0.22\textwidth]{att2_seasonal_autumn.png}
    \end{tabular}
    \caption{四季日内负载与光伏曲线}
    \label{fig:eda_season}
  \end{subfigure}

  \caption{数据探索性分析：(a) 典型日、(b) 全年净负载、(c) 全年电价、(d) 四季曲线}
  \label{fig:eda}
\end{figure}

由图~\ref{fig:eda} 可读出：
\begin{enumerate}[label=\arabic*.)]
  \item 典型日（图 a）电价呈“夜间低谷、早晚双峰、正午次低谷”形态，全天最低 $0.3713$ 元/kWh（05:40）、最高 $1.3952$ 元/kWh（20:40）；光伏仅在 $05{:}00$--$19{:}00$ 出力、正午峰值约 $7600$\,kW，负载 $08{:}00$--$17{:}00$ 维持高位。光伏与电价低点重叠于正午、电价高峰在早晚，为储能“低价充电、高价放电”提供了空间。
  \item 全年净负载持续曲线（图 b）显示净负载为正的时段占 $80.7\%$、范围 $[-6602,6688]$\,kW，正负两端均超过接口功率 $5000$\,kW；日电量上负载相对平稳、光伏夏高冬低、净缺口与光伏大致反向。储能既不能完全消纳光伏过剩，也不能补足最大缺口，模型必须同时设置弃光与紧急购电两个变量。
  \item 全年电价（图 c）日内峰谷明显、范围 $0.0076\sim1.7936$ 元/kWh，其中与附件 1 相关性最高的日期为 2025-04-22（相关系数 $0.992$），说明附件 1 足以代表全年电价形态。
  \item 四季曲线（图 d）显示负载四季差异较小但夜间始终偏高，光伏夏季日均 $2673$\,kW、约为冬季 $1681$\,kW 的 $1.6$ 倍，季节差异显著，光伏预测须引入季节特征。
\end{enumerate}

此外，工作日负载峰值 $6195$\,kW、周末 $5367$\,kW，工作日高约 $15\%$，其星期结构由问题二的负载回归处理；附件 3 的 $0{:}00$ 光伏预报白天平均 MAE 约 $311$\,kW，晴天与多云可差一个数量级，详见问题三。

\subsection{数据处理结果小结}

综上，原始数据完整、无缺失、无负值，异常点均为真实波动予以保留，处理工作集中在时间口径对齐与面向模型的特征构造。处理后得到统一 10 分钟尺度的规范数据集：附件 2、4 各 $365\times144=52560$ 条，附件 1 为 144 条，附件 3 对齐至 10 分钟尺度，并构造出问题二所需的 11 个特征。后续问题一基于附件 1 开展确定性线性规划，问题二基于附件 2 开展滚动预测、日前优化与回测，问题三、四基于附件 3、4 开展不确定性分析与全年电价建模。